\documentclass[10pt]{article}
\usepackage[a4paper, margin=0.75in]{geometry}
\usepackage{titlesec}
\usepackage{booktabs}
\usepackage{tabularx}
\usepackage{xcolor}
\usepackage{enumitem}
\usepackage{tgheros}
\usepackage[T1]{fontenc}
\usepackage{array}
\usepackage{textcomp}
\usepackage[hidelinks]{hyperref}
\renewcommand{\familydefault}{\sfdefault}

%% ===== Colour palette =====
\definecolor{primary}{RGB}{15,52,96}
\definecolor{accent}{RGB}{198,124,0}
\definecolor{lightgrey}{gray}{0.88}
\definecolor{darkgrey}{gray}{0.38}

%% ===== Section style =====
\titleformat{\section}{%
  \color{primary}\bfseries\large\raggedright
}{}{0em}{}[\color{lightgrey}{\titlerule[1.5pt]}\vspace{-4pt}]

\titlespacing*{\section}{0pt}{16pt}{6pt}

\setlength{\parindent}{0pt}
\setlength{\parskip}{6pt}
\renewcommand{\arraystretch}{1.25}
\hypersetup{colorlinks=true, urlcolor=primary}

\newcommand{\spacer}{\vspace{10pt}}
\newcommand{\bigspacer}{\vspace{16pt}}

%% =====================================================
\begin{document}

%% ----- Sender block -----
\begin{flushright}
  {\small
  \textbf{Olajide Olaniyan}\\
  CEng Software Engineer, IBM\\
  Dublin, Ireland\\[2pt]
  +353 89 982 7429\\
  \href{mailto:jidemobell@gmail.com}{jidemobell@gmail.com}\\
  \href{https://linkedin.com/in/olajide-olaniyan}{linkedin.com/in/olajide-olaniyan}\\[4pt]
  {\color{darkgrey}May 2026}
  }
\end{flushright}

\spacer

%% ----- Recipient block -----
{\small
\textbf{[Colleague's Name]}\\
Newcross and Pan Ocean Group of Companies\\
Lagos, Nigeria
}

\bigspacer

%% ----- Subject -----
\begin{center}
  {\color{primary}\large\bfseries Expression of Interest ---
  Digital Transformation Conversation}
\end{center}

\spacer

Dear [Colleague's Name],

Thank you for reaching out about the group's plans around digital
transformation. As someone who spent eleven years building and leading Pan
Ocean's IT infrastructure from the ground up, this is not a conversation
I could approach lightly.

Since leaving in 2021 I have deliberately shaped my career to bridge the
gap between enterprise oil and gas operations and modern software
engineering. My M.Sc.\ in Big Data Analytics and Artificial Intelligence
from Atlantic Technological University (Ireland), combined with my current
role as a CEng Software Engineer at IBM working on Cloud Pak for AIOps ---
event-driven, AI-powered operations on Kubernetes --- has given me the
toolkit this kind of programme calls for.

I want to be honest about what I am offering. I do not assume to know what
is already in motion inside the group. RADA-AMS, infrastructure
modernisation, and other internal initiatives almost certainly have
momentum I cannot see from outside. What I am offering is a perspective
paper, an external view of where the industry sits in 2026, and the
relationships and execution capacity to add to what is already being done.

\bigspacer

%% ----------------------------------------------------------
\section{What I Bring}
%% ----------------------------------------------------------

\noindent\colorbox{primary!8}{%
\begin{minipage}{\dimexpr\textwidth-2\fboxsep}
  \vspace{6pt}
  \textbf{\color{primary}The Dual Lens.}\quad
  I left Pan Ocean to gain skills Nigeria couldn't give me yet. I now have
  them. I am returning with a clear intention: to help build the future of a
  group whose foundation I helped lay.
  \vspace{6pt}
\end{minipage}}

\spacer

\small
\begin{tabularx}{\columnwidth}{@{}XX@{}}
\toprule
\textbf{Inside the Group (2010--2021)} & \textbf{Since the Group (2021--2026)} \\
\midrule
Built Pan Ocean's data centre and WAN infrastructure spanning Lagos HQ,
Ogharefe (OML-98), Owa Aladinma (OML-147), and OOGP field sites
  & CEng Software Engineer at IBM --- Cloud Pak for AIOps:
    AI-powered operations using Java microservices, Kafka, and distributed
    systems on Kubernetes/OpenShift \\
Led migration to Microsoft Azure hybrid cloud
  & M.Sc.\ Big Data Analytics and Artificial Intelligence,
    Atlantic Technological University, Ireland \\
Managed 500+ users across offshore and onshore operations; final escalation
point for the group's most complex technical issues
  & Founding Engineer at Klem AI: product from concept to App Store launch
    (Python, React Native, TensorFlow) \\
Implemented Nagios monitoring infrastructure: 40\% cost reduction across
multi-site operations
  & Associate Operations Engineer, Activision Blizzard: 99.9\%+
    availability at enterprise scale \\
Established security, firewall, and disaster recovery protocols for pipeline
and gas plant infrastructure
  & AI knowledge systems, RAG architecture, Python / LLM integration,
    Kubernetes, Kafka \\
\bottomrule
\end{tabularx}
\normalsize

\spacer
I know Pan Ocean's infrastructure, culture, and operational challenges from
the inside. I know where the friction lives in the business because I lived
it. At IBM I work on the systems that replace that friction with intelligence.
I speak both languages.

\bigspacer

%% ----------------------------------------------------------
\section{Why This Conversation, Why Now}
%% ----------------------------------------------------------

Several factors have converged that make digital capability investment more
pressing than at any point in the group's history:

\spacer
\small
\begin{tabularx}{\columnwidth}{@{}lX@{}}
\toprule
\textbf{Factor} & \textbf{Why It Matters Now} \\
\midrule
RADA-AMS (October 2025)
  & The group publicly committed to energy and innovation leadership. That
    commitment needs an execution engine. \\
NUPRC April 2026 GHG Directive
  & Measurement-based methane reporting is now mandatory. OOGP is directly
    in scope. Manual compliance is not operationally sustainable at scale. \\
PIA 2021 HCDT Obligations
  & 3\% of opex per entity must be tracked, reported, and audited. Across
    three entities, a significant compliance burden if done manually. \\
Three-Entity Group Structure
  & Pan Ocean, Newcross Petroleum, and Newcross Exploration and Production
    operating with shared infrastructure and separate legal reporting is a
    digital integration problem of real complexity. \\
OML-147 PSC Live Reporting
  & Active PSC obligations to NNPC require continuous and accurate
    production and cost recovery data. \\
Fifty Years of OML-98 Operational Data
  & Sitting largely inaccessible. Document intelligence and data engineering
    can unlock it before another decade passes. \\
\bottomrule
\end{tabularx}
\normalsize

\bigspacer

%% ----------------------------------------------------------
\section{What I Have Prepared}
%% ----------------------------------------------------------

Attached to this letter are two papers I have prepared in advance of any
conversation:

\spacer
\small
\begin{tabularx}{\columnwidth}{@{}lX@{}}
\toprule
\textbf{Document} & \textbf{Contents} \\
\midrule
\textit{Digital Transformation --- A Perspective Paper}
  & An external industry view of digital transformation in oil and gas in
    2026, framed around eight capability themes (connected operations,
    industrial AI, data fabric, ESG digitisation, workforce digitisation,
    OT security, generative AI for institutional memory, hybrid cloud and
    edge). Adapted to the group's specific context, the Nigerian regulatory
    environment, and the realities of operating at OOGP, OML-98, OML-147,
    and along the AEPP corridor. \\
\textit{AI and Data Roadmap}
  & A concrete, asset-specific catalogue of AI and data engineering
    initiatives, organised by data and infrastructure readiness rather than
    by calendar dates. Every initiative is grounded in actual group assets
    and Nigerian operational constraints. No research projects.
    No Silicon Valley assumptions. \\
\bottomrule
\end{tabularx}
\normalsize

\spacer
These are discussion papers. They are not a delivery plan and do not assume
visibility of work already in flight. The best outcome of sharing them is
a conversation where I can hear what leadership is actually prioritising
and refine the thinking accordingly.

\bigspacer

%% ----------------------------------------------------------
\section{How I Am Thinking About Engagement}
%% ----------------------------------------------------------

I am open to whatever engagement structure makes sense given the group's
existing plans and team:

\spacer
\small
\begin{itemize}[leftmargin=1.8em, itemsep=4pt]
  \item A \textbf{technical leadership role} --- architecture, standards,
        team-building, and strategic direction alongside delivery, if the
        group is establishing a new function
  \item A \textbf{consulting arrangement} --- structured around specific
        capability areas, working alongside an existing team
  \item An \textbf{advisory engagement} --- contributing perspective,
        architecture review, and external industry context to ongoing work
  \item A \textbf{transitional structure} --- beginning on a defined scope,
        with mutual flexibility to expand as the work and the team take shape
\end{itemize}
\normalsize

\spacer
I will have my Irish Stamp 4 immigration status by mid-2026, which gives me
full flexibility to work formally with additional companies. Any remuneration
would be structured to a Nigerian account.

\bigspacer

%% ----------------------------------------------------------
\section{Questions Worth Discussing}
%% ----------------------------------------------------------

When we speak, these are the things most worth clarifying first --- and what
I would most want to learn from leadership:

\spacer
\small
\begin{enumerate}[leftmargin=1.8em, itemsep=4pt]
  \item What are the \textbf{top two or three business problems} the group
        most wants visibly solved over the next year?
  \item What is \textbf{already in motion} that I should understand before
        proposing anything --- inside RADA-AMS, infrastructure modernisation,
        or other initiatives?
  \item What is the \textbf{governance and reporting structure} for digital
        capability work across the three entities?
  \item Where is the \textbf{commercialisation conversation} today --- is
        external productisation on the table or strictly internal for now?
  \item What is the \textbf{hiring posture} --- internal only, or open to
        diaspora engineers and external technical hires?
  \item What does \textbf{success look like} at twelve months? At thirty-six?
\end{enumerate}
\normalsize

\spacer
I would welcome a conversation with you and the relevant stakeholders. The
two papers give the substance; the conversation is where I can listen and
adapt to what the group is actually building.

\bigspacer

Warm regards,

\vspace{18pt}

\textbf{Olajide Olaniyan}\\
{\small CEng Software Engineer, IBM}\\
{\small +353 89 982 7429\quad|\quad
\href{mailto:jidemobell@gmail.com}{jidemobell@gmail.com}\quad|\quad
\href{https://linkedin.com/in/olajide-olaniyan}{linkedin.com/in/olajide-olaniyan}}

\spacer
{\small\color{darkgrey}
\textit{Attachments: Digital Transformation --- A Perspective Paper (May 2026)
$\cdot$ AI and Data Roadmap (May 2026) $\cdot$ CV}
}

\end{document}
